% Appendix: prompt templates and evaluation items

\subsection*{Appendix A: Sentiment Classification Prompt Templates}
\label{app:sentiment-prompts}

\paragraph{Zero-shot template.}
\begin{verbatim}
You are a sentiment classifier for product reviews.
Some reviews use sarcasm, irony, or mixed language. Focus on the reviewer's
true overall intent, not surface-level word choice.
Classify the review as exactly one word: Positive or Negative.

Review: "{review}"

Label:
\end{verbatim}

\paragraph{One-shot template.}
\begin{verbatim}
You are a sentiment classifier for product reviews.
Some reviews use sarcasm, irony, or mixed language. Focus on the reviewer's
true overall intent.
Classify the review as exactly one word: Positive or Negative.

Example:
Review: "Congratulations to the manufacturer for making something that breaks
in record time. Truly an achievement."
Label: Negative

Now classify:
Review: "{review}"
Label:
\end{verbatim}

\paragraph{Few-shot template.}
\begin{verbatim}
You are a sentiment classifier for product reviews.
Some reviews use sarcasm, irony, or mixed language. Focus on the reviewer's
true overall intent.
Classify the review as exactly one word: Positive or Negative.

Examples:

Review: "Oh great, another charger that works for exactly one month.
Money well spent."
Label: Negative

Review: "I was skeptical after the bad reviews but honestly this thing
exceeded every expectation I had."
Label: Positive

Review: "It's not pretty and the instructions were useless, but once it's
set up it runs like a dream."
Label: Positive

Review: "Sure, it looks nice on the shelf. Too bad it doesn't actually do
what it's supposed to."
Label: Negative

Review: "I bought three more for my family — that should tell you everything."
Label: Positive

Now classify:
Review: "{review}"
Label:
\end{verbatim}

\newpage
\subsection*{Appendix B: Chain-of-Thought Prompt Templates}
\label{app:cot-prompts}

\paragraph{Direct-answer template.}
\begin{verbatim}
You are solving a short arithmetic word problem.
Read the problem carefully and provide only the final numeric answer.
Do not show your reasoning.

Problem: {question}

Answer (just the number):
\end{verbatim}

\paragraph{Chain-of-thought template.}
\begin{verbatim}
You are solving a short arithmetic word problem.
First, think through the problem step by step.
Then, on the last line, provide only the final numeric answer.

Problem: {question}

Let's reason step by step, then give the final answer.
\end{verbatim}

\newpage
\subsection*{Appendix C: Sentiment Evaluation Set}
\label{app:sentiment-data}

The ten hand-crafted product reviews used in Problem~1, with gold labels:
\begin{enumerate}
  \item ``After the update it became the most expensive paperweight in my house.'' (Negative)
  \item ``The replacement works perfectly.'' (Negative)
  \item ``Thanks to this thing I haven't slept properly in a week. I just can't put it down.'' (Positive)
  \item ``It finally convinced me to switch to the competitor.'' (Negative)
  \item ``It didn't catch fire. Exceeds expectations for this brand.'' (Negative)
  \item ``I left a one-star review six months ago. I was wrong.'' (Positive)
  \item ``This ruined every other speaker I own. Nothing else sounds acceptable anymore.'' (Positive)
  \item ``Fantastic product if you don't mind the smell of burning plastic every time you turn it on.'' (Negative)
  \item ``I've been giving these away to everyone I know. Perfect gift for people I don't particularly like.'' (Negative)
  \item ``Bought this four years ago. Still here. Still works. Nothing else to say.'' (Positive)
\end{enumerate}

\newpage
\subsection*{Appendix D: Arithmetic Problems for Chain-of-Thought Evaluation}
\label{app:cot-data}

The arithmetic word problems used in Problem~2, with correct numeric answers:
\begin{enumerate}
  \item A store sells notebooks for \$4 each. If you buy 3 or more, you get a 25\% discount on the entire purchase. How much do 5 notebooks cost? (Answer: 15)
  \item A train travels at 60 mph for 2.5 hours, then at 40 mph for 1.5 hours. What is the total distance traveled in miles? (Answer: 210)
  \item Tom is twice as old as Jerry. In 5 years, Tom will be 35. How old is Jerry now? (Answer: 15)
  \item A recipe calls for $\tfrac{2}{3}$ cup of sugar for 4 servings. How many cups of sugar do you need for 18 servings? (Answer: 3)
  \item A pool fills at 12 gallons per minute and drains at 4 gallons per minute simultaneously. How many minutes does it take to accumulate 120 gallons? (Answer: 15)
  \item Three friends split an \$84 dinner bill. They add a 20\% tip before splitting. How much does each person pay in dollars? (Answer: 33.6)
  \item A shirt originally costs \$60. It is marked down 30\%, then an additional 10\% off the sale price. What is the final price? (Answer: 37.8)
  \item You have 100 feet of fencing to make a rectangular garden. If the length is 10 feet more than the width, what is the area in square feet? (Answer: 600)
  \item A car uses 3 gallons of gas every 90 miles. Gas costs \$4 per gallon. How much does gas cost for a 450-mile trip? (Answer: 60)
  \item A baker makes 48 cupcakes. She sells one third on Monday, then half of what remains on Tuesday. How many are left? (Answer: 16)
  \item Two runners start from the same point running in opposite directions. One runs 7 mph and the other 5 mph. After how many hours are they 36 miles apart? (Answer: 3)
  \item A company's profit was \$200{,}000 last year. It grew 15\% this year but they paid 40\% in taxes on this year's profit. What is the after-tax profit this year? (Answer: 138000)
\end{enumerate}

\newpage
\subsection*{Appendix E: Prompt Variants for Sensitivity Analysis}
\label{app:sensitivity}

For Problem~3 (prompt sensitivity), we used the Problem~1 few-shot sentiment prompt as the baseline and produced four controlled perturbations on the same five labeled examples and task description.

\begin{itemize}
  \item \textbf{Baseline.} The unmodified few-shot template from Appendix~A (order: Neg, Pos, Pos, Neg, Pos).
  \item \textbf{Reorder.} Same five examples and wording; only the order of examples changed to positive-first (Pos, Pos, Pos, Neg, Neg).
  \item \textbf{Rephrase.} Same example order as baseline; the task line was shortened to ``Classify the sentiment as Positive or Negative'' (removed ``exactly one word'').
  \item \textbf{Add constraint.} Baseline plus one sentence: ``Do not explain; output only the label.''
  \item \textbf{Remove constraint.} Baseline with ``exactly one word'' removed from the instruction (``Classify the review as Positive or Negative'').
\end{itemize}

The full text of each variant follows the same structure as the few-shot template in Appendix~A (task description, five examples, ``Now classify:'', then the query review and ``Label:''). Implementations are in \texttt{src/incontext\_learning.py} under \texttt{SENSITIVITY\_VARIANTS}.
